% templates/transition_exhibit.tex.mustache
% AUTO-GENERATED transition exhibit section
% Rendered by scripts/build_campaign_report.jl
% DO NOT EDIT MANUALLY

\section{Transition Event Mapping}

Transition plots are the primary dynamical fingerprint for regime transitions in the stable boundary layer. This section provides a three-panel bridge between continuation stability signatures and reduced-state trajectory structure for campaign \texttt{ BLLAST }.

\begin{figure}[htbp]
    \centering
    \usepgfplotslibrary{groupplots}
    \tikzsetnextfilename{transition-panel-bllast}
    \begin{tikzpicture}
        \begin{groupplot}[
            group style={group size=1 by 3, vertical sep=14pt},
            width=0.85\textwidth,
            height=5.2cm,
            grid=both,
            tick label style={font=\small},
            label style={font=\small},
            title style={font=\small\bfseries, align=center},
            axis background/.style={fill=gray!2}
        ]

        \nextgroupplot[
            title={Panel A: Stability Branch (S-Curve Proxy)},
            xlabel={$\gamma$},
            ylabel={$\alpha=\max\Re(\lambda)$}
        ]
            \addplot[
                color=blue!70!black,
                line width=1.2pt,
                mark=* ,
                mark size=1.4pt,
                mark options={fill=blue!70!black}
            ] table[
                col sep=comma,
                x=gamma,
                y=max_real_eig
            ] {../../data/outputs/transition_panel_a_bllast.csv};

        \nextgroupplot[
            title={Panel B: State-Space Projection},
            xlabel={$z_1$},
            ylabel={$z_2$}
        ]
            \addplot[
                only marks,
                scatter,
                scatter src=explicit,
                mark size=1.8pt
            ] table[
                col sep=comma,
                x=z1,
                y=z2,
                meta=z3
            ] {../../data/outputs/transition_panel_b_bllast.csv};

        \nextgroupplot[
            title={Panel C: Reconstructed $\mathrm{Ri}_g(z,\gamma)$ Structure},
            xlabel={$\gamma$},
            ylabel={Height $z$ [m]},
            ymin=2.0, ymax=60.0,
            axis lines=left,
            point meta min=-0.5,
            point meta max=2.5,
            colormap/viridis,
            colorbar,
            colorbar style={title={$\mathrm{Ri}_g$}},
            legend style={at={(0.98,0.05)}, anchor=south east, font=\footnotesize},
            axis background/.style={fill=gray!2}
        ]
            \addplot[
                matrix plot*,
                mesh/cols=5,
                point meta=explicit
            ] table[
                col sep=comma,
                x=gamma,
                y=height_z,
                meta=ri_g
            ] {../../data/outputs/transition_panel_c_profile_bllast.csv};

            \addplot[
                only marks,
                mark=* ,
                mark size=0.5pt,
                color=red!80!black,
                forget plot
            ] table[
                col sep=comma,
                x=gamma,
                y=height_z,
                row predicate/.code={
                    \pgfmathparse{abs(\thisrow{ri_g}-0.25) < 0.02 ? 1 : 0}
                    \ifnum\pgfmathresult=1\relax\else\pgfplotstableuserowfalse\fi
                }
            ] {../../data/outputs/transition_panel_c_profile_bllast.csv};
            \addlegendentry{$\mathrm{Ri}_c \approx 0.25$ band}

            \node[anchor=west, text=teal!90!black, font=\scriptsize\bfseries] at (axis description cs:0.03,0.94) {$\eta_3$ Curvature Inflexion};

        \end{groupplot}
    \end{tikzpicture}
    \caption{Three-panel transition exhibit. Panel A tracks the stability branch and the Hopf-like crossing. Panel B projects the continuation branch in reduced manifold coordinates. Panel C shows the localized $\eta_3$ transition window from campaign trajectories.}
    \label{fig:transition_exhibit_bllast}
\end{figure}

\paragraph{Transition Metrics}
\begin{itemize}
    \item Critical coupling estimate: n/a
    \item Hopf period estimate: 6.20 min
    \item Peak $|\eta_3|$ in transition window: 3.1604
\end{itemize}

\subsection{Structural Cross-Diagnostics: Curvature vs. Local Criticality}

\begin{figure}[htbp]
    \centering
    	ikzsetnextfilename{structural-phase-portrait-bllast}
    \begin{tikzpicture}
        \begin{axis}[
            width=0.74\textwidth,
            height=0.52\textwidth,
            grid=major,
            xlabel={Gradient Richardson Number $\mathrm{Ri}_g$ (lowest mast level)},
            ylabel={Structural Curvature Amplitude $\eta_3$},
            title={Manifold Phase Portrait: BLLAST},
            extra x ticks=0.25,
            extra x tick labels={$\mathrm{Ri}_c$},
            extra x tick style={grid=major, tick label style={above, yshift=4pt}, line width=1.0pt, dashed, color=red},
            legend style={at={(1.04,1.0)}, anchor=north west, draw=none, fill=none}
        ]
            \addplot[
                only marks,
                mark=* ,
                mark size=1.2pt,
                opacity=0.4,
                color=blue!55!magenta
            ] table[
                col sep=comma,
                x={ri_g_low},
                y={eta_3}
            ] {../../data/outputs/structural_correlations_bllast.csv};
            \addlegendentry{$\eta_3$ vs $\mathrm{Ri}_g$}
        \end{axis}
    \end{tikzpicture}
    \caption{Phase portrait mapping continuous column curvature ($\eta_3$) against local gradient stability. The dashed line at $\mathrm{Ri}_c = 0.25$ marks the classical local threshold while the manifold curvature remains dynamically active across the transition.}
    \label{fig:curvature_vs_ri_bllast}
\end{figure}

\paragraph{Structural Metrics}
\begin{itemize}
    \item Low-level Richardson proxy column: \texttt{ ri\_g\_2\_0 }
    \item Best lag correlation between $d\eta_3/dt$ and $d\mathrm{Ri}_g/dt$: $R_\tau=0.248$ at lag $\tau=0.00\,\mathrm{h}$ (0 samples)
    \item Segmented mutual information (bits): $I(\eta_3;\mathrm{Ri}_g\leq \mathrm{Ri}_c)=1.453$, $I(\eta_3;\mathrm{Ri}_g>\mathrm{Ri}_c)=n/a$
\end{itemize}

